Only one assignment is observed from each cluster schedule.  Consequently \(P_0\) and the product coupling \(\bar P_0=\bigotimes_zP_{0,z}\) have the same marginal law at the realized assignment, proving \(Q^B_{\bar P_0,p}=Q^B_{P_0,p}\) by \(\texttt{def:observed-law-kernels}\).  The assumptions \(\texttt{ass:iid-schedules}\), \(\texttt{ass:isolated-clusters}\), \(\texttt{ass:bernoulli-label-iid}\), \(\texttt{ass:bernoulli-label-schedule-independence}\), and \(\texttt{ass:bernoulli-units}\) make the cluster records independent and give the exact-hit probabilities in \(\texttt{lem:exact-hit-factorization}\).

Fix \(h\in\mathbb R^A\), without imposing a zero-sum normalization.  For \(z\in\mathcal S_k\), put \(V_z=W(z)-\mu_z\).  Because the lifted outcome domain gives \(0\leq W(z)\leq1\), \(|V_z|\leq1\), so the bounded-support exponential tilt in \(\texttt{def:least-favourable-path}\) is well defined for every finite tilt.  Differentiation under the bounded integral gives
\[
\left.\frac{d}{dt}\mathbb E_{P_{t,h,z}}W(z)\right|_{t=0}
=\frac{h_k}{\tau_k^2}\operatorname{Var}_{P_0}\{W(z)\}. \tag{17}
\]
Division is valid by \(\texttt{ass:nondegenerate-active-slices}\).  Averaging (17) over \(\mathcal S_k\) and using the definition of \(\tau_k^2\) yields
\[
\left.\frac d{dt}U_k(\bar P_{t,h})\right|_{t=0}
=\frac{h_k}{M_k\tau_k^2}\sum_{z\in\mathcal S_k}
  \operatorname{Var}_{P_0}\{W(z)\}=h_k. \tag{18}
\]

The observed score in coordinate \(k\) is \(\dot\ell_k^B\).  It is centered and bounded.  Distinct count slices are disjoint, and \(\texttt{lem:exact-hit-factorization}\) makes the conditional assignment uniform within a hit slice.  Therefore
\[
\mathbb E_0(\dot\ell_j^B\dot\ell_k^B)
=\mathbf1\{j=k\}\frac{q_k}{M_k\tau_k^4}
  \sum_{z\in\mathcal S_k}\operatorname{Var}_{P_0}\{W(z)\}
=\mathbf1\{j=k\}\frac{q_k}{\tau_k^2}. \tag{19}
\]
This is \(I_B\), including the common-shift score \(\sum_{k\in A}\dot\ell_k^B\).

Bounded support also gives a uniform cumulant expansion.  Under every tilted law the third derivative of the log moment-generating function is the third centered moment of a variable in an interval of length one.  Hence, for \(|a|\) in any fixed compact interval,
\[
\Lambda_{z,k}(a)=\frac{a^2}{2}\operatorname{Var}_{P_0}\{W(z)\}+R_{z,k}(a),
\qquad |R_{z,k}(a)|\leq C_0|a|^3, \tag{20}
\]
with one finite \(C_0\) because there are finitely many assignments.  Substituting \(a=C^{-1/2}h_k/\tau_k^2\), summing the one-cluster likelihood contributions, and using the bounded law of large numbers for assignment-cell frequencies gives
\[
\log\frac{d(Q^B_{\bar P_{C,h},p})^{\otimes C}}
         {d(Q^B_{P_0,p})^{\otimes C}}
=h^{\prime}\Delta_C^B-\frac12h^{\prime}I_Bh+o_P(1). \tag{21}
\]
The total third-order remainder is \(O(C^{-1/2}\max_k|h_k|^3)\) uniformly on compact sets of representatives.  The summands of \(\Delta_C^B\) are independent, centered, bounded vectors with covariance (19).  For every fixed vector \(u\), their characteristic functions satisfy
\[
\mathbb E_0\exp\{i u^{\prime}\dot\ell_A^B/\sqrt C\}
=1-\frac{u^{\prime}I_Bu}{2C}+O(C^{-3/2}\lVert u\rVert_2^3). \tag{22}
\]
Raising (22) to the \(C\)-th power proves \(\Delta_C^B\Rightarrow N(0,I_B)\), thereby deriving rather than assuming the LAN limit.

The asymptotic linear representation from \(\texttt{thm:efficient-exact-hit-bernoulli}\) has influence coordinate \(\phi_k^B\).  Directly from the definitions and (19),
\[
\mathbb E_0(\phi_j^B\dot\ell_k^B)=\mathbf1\{j=k\}. \tag{23}
\]
The joint characteristic-function calculation for \((\phi_A^B,\dot\ell_A^B)\), followed by exponential tilting through (21), gives under \(\bar P_{C,h}\)
\[
\sqrt C\{\widetilde U_{C,A}-U_A(P_0)\}
\Rightarrow N\{h,\Sigma_{B,A}\}. \tag{24}
\]

For the decision lower bound, take any finite collection of representatives satisfying \(\lVert[h]\rVert_A\leq M\) and any finite prior on them.  Equations (21)--(24) verify the LAN and loss-convergence premises of \(\texttt{lem:lan-finite-prior-lower-bound}\); bounded likelihood scores give uniform integrability.  Hence every rule sequence has limiting maximum risk at least that finite prior's Gaussian Bayes risk.  To retain the unknown common location, tensor a finite quotient prior with a uniform grid of shifts \(c\mathbf1_A\), \(-T\leq c\leq T\).  Gaussian total-variation continuity makes the mesh error vanish, while the boundary-to-volume ratio of the shift grid vanishes as \(T\to\infty\).  Averaging the rule over this grid yields a location-invariant rule.  Finite parameter grids are dense on the compact quotient ball because Gaussian risks are uniformly continuous there.  Taking suprema over the finite priors and invoking \(\texttt{lem:gaussian-quotient-reduction}\) gives precisely \(G_{A,M}(\tau_A^2,q_A)\), not a game with the common location known.  The fixed welfare gap below the active face makes every inactive action asymptotically dominated.  The same quotient lemma gives \(G_{A,M}\uparrow G_A\), and \(\texttt{thm:nuisance-deletion-bernoulli}\) transfers the lower bound to full-record rules.

The early finite check agrees with these formulas.  For \(n=4\), copy one independent \(\operatorname{Bernoulli}(1/2)\) target bit to all four unit outcomes at each target assignment.  Then \(W\) is that bit, \(\tau_k^2=1/4\), information is \(4q_k\), and calibrated influence variance is \(1/(4q_k)\) on both target slices.  On the four count-one assignments, instead take legal copied-bit means \((4/5,4/5,1/5,1/5)\).  Then \(U_1=1/2\), every assignment variance is \(4/25\), \(\tau_1^2=4/25\), information is \(25q_1/4\), and influence variance is \(4/(25q_1)\).  Thus assignment heterogeneity changes \(\tau_k^2\) but not the exact-hit multiplier.  For two tied policies, writing \(V=\tau_1^2/q_1+\tau_2^2/q_2\), the local pair \(h=(d/2,-d/2)\) has wrong-selection regret \(d\Phi(-d/\sqrt V)\); maximizing over \(d\geq0\) gives \(\kappa\sqrt V\), exactly the pooled Gaussian regret in \(\texttt{thm:two-policy-minimax}\).  Finally, the empty-cell fallback has finite-sample bias
\[
\frac1{M_k}\sum_{z\in\mathcal S_k}(1/2-\mu_z)
\left(1-\frac{q_k}{M_k}\right)^C,
\]
which is exponentially small because \(q_k/M_k>0\); no finite-sample unbiasedness is used.